2-1 Insertion sort on small arrays in merge sort\\
Although merge sort runs in $\theta(nlgn)$ worst-case time and insertion sort runs 
in $\theta(n^2)$ worst-case , the constant factors in insertion sort can make it faster 
in practice for small problem sizes on many machines. Thus it makes sense to 
coarsen the leaves of the recursion by using insertion sort within merge sort when 
subproblems become sufûciently small. Consider a modiûcation to merge sort in 
which n/k sublists of length k are sorted using insertions sort and then merged using
the standard merging mechanism, where k is a vlue to be determined. \\

a. Show that insertion sort can sort the n/k sublists, each of length k, in $\theta(nk)$
worst-case time. \\

\begin{verbatim}
############ CODE FOR THIS SECTION #############
from xmlrpc.client import boolean

def insertionSort(nums: list[int]) -> list[int]:
  n: int = len(nums)

  res: list[int] = []
  
  for num in nums:
    res.append(num)

  if n == 0:
    return []
  if n == 1:
    return [nums[0]]
  
  i: int = 1

  while (i < n):
    j: int = i-1
    while (res[j] > res[j+1] and j >= 0):
      temp: int = res[j+1]
      res[j+1] = res[j]
      res[j] = temp
      j = j-1
    i += 1
  
  return res

def mergeSort(nums: list[int], l: int, r: int, k: int) -> list[int]:

  def merge(nums1: list[int], nums2: list[int]) -> list[int]:
    # nums1
    i: int = 0
    n: int = len(nums1)

    # nums2
    j: int = 0
    m: int = len(nums2)

    res: list[int] = []

    while (i < n and j < m):
      if (nums1[i] <= nums2[j]):
        res.append(nums1[i])
        i += 1
      else:
        res.append(nums2[j])
        j += 1
    
    if i >= n:
      res.extend(nums2[j:m])
    else:
      res.extend(nums1[i:n])
    
    return res
  
  if not nums or l > r:
    return []
  if l == r:
    return [nums[l]]
  if len(nums) <= k:
    return insertionSort(nums)
  
  m: int = (l+r) // 2
  leftNums: list[int] = mergeSort(nums, l, m, k)
  rightNums: list[int] = mergeSort(nums, m+1, r, k)
  return merge(leftNums, rightNums)

array: list[int] = [100, 43, 54, 23, 67, 1, 9, 413, 456, 21, 46, 80, 70]
sortedArray: list[int] = mergeSort(array, 0, len(array)-1, 3)

print(array)
print(sortedArray)

\end{verbatim}

So given n/k sublists each of length k I need to show the worst case will be $\theta(nk)$
let x = n/k which is the number of sublists then we would be calling the sorting algorithm x times 
so the worst case run time would be x*(worst case run time of insertion sort on array of len k)

\begin{verbatim}
Examine:

def insertionSort(nums: list[int]) -> list[int]:
  n: int = len(nums)                      (1)
  i: int = 1                              (2)

  while (i < n):                          (3)
    j: int = i-1                          (4)
    while (res[j] > res[j+1] and j >= 0): (5)
      temp: int = res[j+1]                (6)
      res[j+1] = res[j]                   (7)
      res[j] = temp                       (8)
      j = j-1                             (9)
    i += 1                                (10)
  
  return nums                             (11)
\end{verbatim}

\begin{proof}
(1) (costs $c_1$, times 1)

(2) (costs $c_2$, times 1)

(3) (costs $c_3$, times k)

(4) (costs $c_4$, times k-1)

(5) (costs $c_5$, times ($\sum_{i=1}^{k-1} (i+1$)))

(6-9) (costs $c_6$, $c_7$, $c_8$, $c_9$ |  times ($\sum_{i=1}^{k-1} i$))

(10) (costs $c_{10}$, times k-1)

(11) (costs $c_{11}$, times 1)
\end{proof}

We know by Gauss sum that: $\sum_{i=1}^{m} i = \frac{m(m+1)}{2}$

So we how have as the summation of all the lines for inserstion sort:\\
\begin{align*}
&c_1 + c_2 + kc_3 + (k-1)c_4 + \left(\dfrac{(k-1)k}{2}+k-1\right)c_5 
+ \dfrac{(k-1)k}{2}\left(c_6 + c_7 + c_8 + c_9\right) + (k-1)c_{10} + c_{ll} \\[6pt]
\implies\ &1\left(c_1 + c_2 + c_{ll}\right) + (k-1)\left(c_4 + c_{10}\right) + kc_3 
+ \dfrac{(k-1)k}{2}\left(c_6 + c_7 + c_8 + c_9\right) + \dfrac{(k-1)k}{2}c_5 + kc_5 - c_5 \\
\implies &1(c_1 + c_2 + c_{11} - c_4 - c_5 - c_{10}) + k(c_3 + c_4 + c_5 + c_{10}) + 
\dfrac{(k-1)k}{2}(c_6 + c_7 + c_8 + c_9 + c_5) \\
\implies &c_1 + c_2k + (k^2 - k)(c_3) \\
\implies &c_1 + (c_2-c_3)k + c_3k^2 \\
\end{align*}
so insertion sort costs $O(k^2)$ per sub list and there are n/k sublists so the total cost is $O(k^2 * n/k) = O(nk)$ \\

b. Show how to merge the sublists in $\theta(nlg(n/k))$ worst-case time. \\
if you were to use merge sort for the n/k sub lists then you could get a runtime of $\theta(nlg(n/k))$ \\
so you have n/k sorted sublists you then use merge sort so you are callin merge sort on n/k items \\
and we know the runtime of merge sort in 

c. Given that the modified algorithm runs in $\theta(nk + nlg(n/k))$ worst-case time, 
what is the largest value of k as a function of n for which the modified algorithm
has the same running time as standard merge sort, in terms of $\theta-notation$? \\

d. How should you choose k in practice? \\

2-2 Correctness of bubblesort \\
Bubblesort is a popular, but inefûcient, sorting algorithm. It works by repeatedly 
swapping adjacent elements that are out of order. The procedure BUBBLESORT 
sorts array A[i:n].
\begin{verbatim}
def bubblesort(array: list[int], n: int):
    for i in range(n):
        for j in range(n, i, -1):
            if (array[j] < array[j - 1]):
                temp: int = array[j]
                array[j] = array[j-1]
                array[j-1] = temp
\end{verbatim} 

a. Let $A'$ denote the array A after bubblesort(A,n) is executed. To prove that bubblesort is correct, you need to prove that it terminates and that 
$A'[1] \leq A'[2] \leq .... \leq A'[n]$. (2.5) \\
In order to show that bubblesort actually sorts, what else do you need to prove? \\

b. State precisely a loop invariant for the for loop in lines 2-4, and prove that this 
loop invariant holds. Your proof should use the structure of the loop-invariant 
proof presented in this chapter. \\

c. Using the termination condition of the loop invariant proved in part (b), state 
a loop invariant for the for loop in lines 1-4 that allows you to prove inequal- 
ity (2.5). Your proof should use the structure of the loop-invariant proof pre- 
sented in this chapter. \\

d. What is the worst-case running time of BUBBLESORT? How does it compare 
with the running time of INSERTION-SORT? \\

2-3 Correctness of Horner’s rule  \\
You are given the coefficents $a_0, a_1, a_2,...,a_n$ of a polynomial \\
P(x) = $\sum_{k=0}^{n}a_kx^k$ = $a_0 + a_1x + a_2x^2 + ... + a_{n-1}x^{n-1} + a_nx^n$, \\
and you want to evaluate this polynomial for a given value of x. Horner’s rule 
says to evaluate the polynomial according to this parenthesization: \\
$P(x) = a_0 + x(a_1 + x(a_2 + ... + x(a_{n-1} + xa_n)...))$ \\
The procedure HORNER implements Horner’s rule to evaluate 
P(x), given the coefficents $a_0,a_1,a_2,...,a_n$ in an array A[0:n] and the value of x. 
\begin{verbatim}
def horner(array, n, x):
    p = 0
    for i in range(n-1, -1, -1):
        p = a[i] + x*p
    return p
\end{verbatim}
a. In terms of $\theta-notation$, what is the run time of this procedure? \\

b. Write code to implement the naive polynomial-evaluation algorithm that 
computes each term of the polynomial from scratch. What is the running time 
of this algorithm? How does it compare with HORNER? \\

c. Consider the following loop invariant for the procedure horner:
    At the start of each iteration of the for loopof line 2-3,
    $p = \sum_{k=0}^{n-(i+1)}A[k+i+1]*x^k$. \\
Interpret a summation with no terms as equaling 0. Following the structure 
of the loop-invariant proof presented in this chapter, use this loop invariant to 
show that, at termination, \\
 $p = \sum_{k=0}^{n}A[k]*x^k$. \\

2-4 Inversions\\
Let A[1:n] be an array of n distinct numbers. If $i < j$ and $A[i] > A[j]$, then the pair (i,j)
is called an inversion of A. \\
a. List the five inversions of the array $<2,3,8,6,1>$.\\ 

b. What array with elements form the set ${1,2,...,n}$ has the most inversions? \\
How many does it have? \\

c. What is the relationship between the running time of insertion sort and the 
number of inversions in the input array? Justify your answer. \\

d. Give an algorithm that determines the number of inversions in any permutation 
on n elements in $\theta(nlgn)$ worst-case time. (Hint: Modify merge sort.) \\
